% Reconstructed text from V5 PDF page 43; non-native source
\par\smallskip\noindent\rule{\textwidth}{0.4pt}\par{\sffamily\footnotesize Original V5 page 43 -- recovered text}\par\begin{flushleft}\footnotesize\ttfamily
\par\smallskip
\par\smallskip
\noindent Diagnostics and Serving Infrastructure for MAS 43\par
\par\smallskip
\par\smallskip
\par\smallskip
\noindent cache hit rate. We propose DRICE — Diagnostics evidence, Retention policy, Infrastructure\par
\noindent assumptions, Cost decomposition, Evidence depth — as a structured evaluation framework for\par
\noindent characterizing systems that span both layers.\par
\par\smallskip
\noindent 8.1 Framework Dimensions\par
\noindent D — Diagnostics evidence. What diagnostic signal does the system consume or produce? We\par
\noindent assess: (a) the evidence type (execution traces, attention scores, tool-call logs, prefill NLL, or\par
\noindent none); (b) the capture mechanism (instrumentation, model introspection, or external logging);\par
\noindent and (c) the evidence availability — whether the signal is available at runtime or only post-hoc.\par
\noindent A system that produces diagnostic evidence as a byproduct of serving (e.g., the policy-driven\par
\noindent runtime’s agent-identity embeddings) receives a higher D score than one that requires a separate\par
\noindent diagnostic pipeline.\par
\par\smallskip
\noindent R — Retention policy. How does the system decide what to keep and what to discard? We\par
\noindent characterize the retention policy along three sub-dimensions: (a) the eviction trigger (memory\par
\noindent pressure, TTL expiry, explicit release, or never); (b) the retention criterion (recency, frequency,\par
\noindent cost-benefit, workflow-awareness, or learned); and (c) the diagnostic awareness — whether the\par
\noindent retention criterion considers the diagnostic value of the retained state. Among surveyed systems,\par
\noindent only the policy-driven runtime incorporates observability signals into eviction decisions; no\par
\noindent system explicitly weights diagnostic evidence retention.\par
\par\smallskip
\noindent I — Infrastructure assumptions. What serving-infrastructure constraints does the system\par
\noindent assume? We distinguish: (a) model homogeneity (single-model, multi-adapter, or multi-model);\par
\noindent (b) memory topology (GPU-only, CPU offload, disaggregated, or tiered); (c) multi-tenancy (ded-\par
\noindent icated, shared-nothing, or shared-pool); and (d) workload pattern (request-level, session-level,\par
\noindent workflow-level, or agent-level). Infrastructure assumptions determine whether a diagnostic\par
\noindent method’s evidence requirements are compatible with the serving system’s state management.\par
\par\smallskip
\noindent C — Cost decomposition. We decompose runtime cost into four components: (a) compute\par
\noindent cost (replays, model inference, signal extraction); (b) memory cost (trace storage, cache over-\par
\noindent head, metadata); (c) latency cost (per-request overhead, diagnostic pipeline delay); and (d) com-\par
\noindent munication cost (inter-node transfers, cache synchronization). This decomposition enables\par
\noindent identifying which cost component dominates and whether diagnostic and serving costs share\par
\noindent the same bottleneck.\par
\par\smallskip
\noindent E — Evidence depth. At what level of verification is the system’s performance characterized?\par
\noindent We use the three-tier coding from Section 3.3: T1 (full-text with section-level evidence), T1b\par
\noindent (full-text without page-level citation), and T2 (abstract-verified). Evidence depth determines\par
\noindent how much confidence can be placed in comparative claims.\par
\par\smallskip
\noindent 8.2 Applying DRICE\par
\noindent Table 12 applies the framework to 21 representative systems spanning the diagnostic-serving\par
\noindent spectrum. Three observations emerge:\par
\noindent  1. The D–R gap. Most diagnostic systems (CHIEF, MASPrism, FAMAS, AgenTracer) have\par
\noindent  strong D scores (they consume rich evidence) but no R dimension (they do not interact\par
\end{flushleft}
